\subsection{Applications of Ontologies}
\label{sec:ontology-applications}
\slidesrc{Ontology development: 77--104}

\subsubsection{Some research from the course's group}
\slidesrc{Ontology development: 77--83}

The slides list research papers from the University of Genoa group (with Viviana
Mascardi), showing that ontologies are an active research topic, often combined with
software agents:
\begin{center}
\small
\begin{tabularx}{\linewidth}{@{}l X@{}}
\toprule
\textbf{year} & \textbf{paper (topic)} \\
\midrule
2010 & \emph{Automatic Ontology Matching via Upper Ontologies: A Systematic Evaluation}:
  using upper ontologies (SUMO-OWL, OpenCyc, DOLCE) as a semantic bridge in ontology
  matching (Section~\ref{sec:ontology-matching}); a ``mixed'' method that combines
  structural matching with upper ontologies improves the recall \\
2012 & \emph{Semantic Web and Declarative Agent Languages and Technologies: Current and
  Future Trends} (position paper) \\
2014 & \emph{CooL-AgentSpeak}: agents that exchange plans and use ontology services
  (ontology matching to find relevant plans of other agents) \\
2017 & \emph{Agent-oriented and ontology-driven digital libraries: the IndianaMAS
  experience}: a multi-agent system that classifies documents of a cultural heritage
  domain formalised by an ontology \\
2020 & \emph{OntoScene, a Logic-based Scene Interpreter}: ontologies and Prolog to
  understand visual scenes, applied to rock art \\
2023 & \emph{Integrating Ontologies and Cognitive Conversational Agents in On2Conv} \\
2024 & \emph{Integrating Procedural Ontologies in VEsNA}: procedural ontologies for BDI
  agents \\
\bottomrule
\end{tabularx}
\end{center}

\subsubsection{Application areas}
\slidesrc{Ontology development: 84--104}

The application part integrates material from a tutorial by Ian Horrocks.

\begin{itemize}
  \item \textbf{e-Science, geoscience}: mediating geological maps globally, integrating
        multi-source geodata, common schemas for water resources. The challenges of an
        ``ontology spectrum'' for geological data interoperability: modelling vs.\
        encoding, multilinguality, flexibility and usefulness, mediation and evolution.
        Resources: geoscience-semantics ontologies, the W3C Geo vocabulary, the GeoNames
        ontology; also semantic enhancement of volunteered geographic information
        (user-generated tags on maps, a PhD project in Genoa).
  \item \textbf{e-Science, bioinformatics}: used for ``in silico'' investigations, which
        relate theory and data, e.g.\ data on plant structure and functionalities
        (the \textbf{Plant Ontology}, PO). Resources: the OBO Foundry
        (\url{http://www.obofoundry.org/}), books on ontologies for bioinformatics and
        anatomy ontologies.
  \item \textbf{Medicine}: building and maintaining terminologies such as
        \textbf{SNOMED}, \textbf{NCI} and \textbf{GALEN}. Examples: the brain anatomy in
        GALEN, where a reasoner checks definitions such as which gyri are bounded by
        which sulci; the \textbf{Disease Ontology} (DOID), e.g.\ \emph{viral infectious
        disease} defined as ``a disease by infectious agent that results in infection,
        has material basis in some Viruses''; the ICPC-2 ontology for primary care; an
        ontological framework for a healthcare research cluster.
  \item \textbf{Organising complex and semi-structured information}: UN-FAO (the
        \emph{geopolitical ontology}, a master reference of country names and codes in
        several languages), NASA (\textbf{SWEET}, Semantic Web for Earth and
        Environmental Terminology), Ordnance Survey, General Motors, Lockheed Martin,
        \ldots; Genoa projects: ontology-driven management of digital rock art objects
        (IndianaMAS), \emph{MUSE: MUltilinguality and SEmantics for the citizens of the
        world}, \emph{Ontologica}: ontologies and natural language for railway
        management.
  \item \textbf{Military, government, legal domains}: DARPA, NSA, NIST, SAIC, the UK
        Ministry of Defence, the US Department of Homeland Security, \ldots; e.g.\
        ontologies for e-government and \textbf{LKIF-Core}, a core ontology of legal
        concepts.
  \item \textbf{The Semantic Web} (Section~\ref{sec:semantic-web}).
\end{itemize}

\begin{nbwhy}
In all these areas the problem is the same as in Section~\ref{sec:what-is-an-ontology}:
many sources and many people describe the same things differently. An ontology gives
them a shared, machine-readable vocabulary; ontology learning helps build and populate
it from text; ontology matching connects it with the ontologies of others.
\end{nbwhy}

\begin{nbsummary}[applications]
\begin{itemize}[nosep]
  \item e-Science: geoscience (data interoperability), bioinformatics (Plant Ontology,
        OBO Foundry).
  \item Medicine: SNOMED, NCI, GALEN; Disease Ontology.
  \item Organising complex information: UN-FAO geopolitical ontology, NASA SWEET,
        cultural heritage, railways.
  \item Military, government, legal (LKIF-Core); the Semantic Web.
\end{itemize}
\end{nbsummary}
